\chapter{Conclusiones}

Los siete incrementos completados confirman la viabilidad de la arquitectura y del enfoque metodológico. El sistema cubre RF-01 a RF-14 de extremo a extremo ---autenticación y arranque controlado, catálogos, envíos, asignación y ciclo de vida, trazabilidad inmutable, administración de usuarios, cambio de contraseña e indicadores operativos--- y queda accesible por una dirección HTTPS pública mediante un procedimiento de despliegue reproducible. La separación SPA/API, las políticas de autorización, el modelo versionado y las pruebas automatizadas proporcionan una base mantenible sin haber ampliado el alcance más allá del núcleo necesario.

La experiencia acumulada muestra el valor de comenzar con un esqueleto transversal y añadir después rebanadas verticales. Los contratos de error, la configuración de EF en pruebas, la sesión y el proxy se validaron antes de que creciese el dominio; sobre esa base, cada regla se pudo contrastar en servicio, HTTP e interfaz.

La lección más útil, sin embargo, procede de las deudas. Cada sprint declaró por escrito lo que no había resuelto ---la sesión en \texttt{localStorage}, la vigencia de los claims de un JWT ya emitido, las ventanas de concurrencia en la doble reserva y en la protección del último administrador--- en lugar de presentar el incremento como completo. El endurecimiento las cerró una por una, y esa continuidad entre lo declarado y lo resuelto es más informativa que un relato sin fricción. El propio proceso lo confirmó de nuevo al final: una revisión posterior al séptimo incremento encontró trece defectos de coherencia, en su mayoría restos que el cambio de mecanismo de sesión había dejado con apariencia de estar activos. Sustituir algo transversal no termina cuando lo nuevo funciona, sino cuando lo viejo deja de aparentar que sigue en uso.

Un intercambio merece quedar enunciado sin adorno: la cookie \texttt{HttpOnly} reduce la exposición del token frente a XSS a cambio de introducir superficie CSRF, que la operación en mismo origen y \texttt{SameSite=Strict} cubren. No es una mejora gratuita, sino un cambio de riesgo elegido con conocimiento de ambos lados.

El trabajo alcanza así su objetivo sin reclamar más de lo demostrado. El entorno desplegado es una demostración reproducible y no un servicio operado: el túnel es temporal, no hay copias de seguridad ni recolección de métricas, y esa distinción se mantiene visible a lo largo de la memoria. Las líneas futuras fuera del TFG podrán incluir optimización de rutas, seguimiento por GPS, facturación o acceso de clientes externos, siempre después de validar la utilidad del núcleo operativo con uso real.
